\section{System Design}
\label{sec:design}

\system selects which Agents' histories deserve protection (O1),
which KV regions within those histories should be retained (O2), and
how much HBM replication may occupy as foreground demand changes (O3). 
Figure~\ref{fig:system-overview} shows our four cooperating components.
\emph{Tracker} estimates remaining workflow work and recovery urgency.
\emph{Router} chooses execution workers and replica destinations.
\emph{KV Master} selects Agent histories and KV regions, evaluates
candidate placements, and manages replica retention and reclamation.
\emph{Adapter} adjusts the HBM replication budget to foreground demand. Rank-local
Connectors execute transfers and report completion to KV Master.

\begin{figure}[t]
 \centering
 \includegraphics[width=\columnwidth]{figures/system-overview.pdf}
 \caption{\system: Four-component architecture. Each workflow node is one
 invocation. Letters identify Agents, and parentheses give invocation
 ordinals, B and C are each invoked twice.
 Each host contains one TP=2 DP worker. Each rank's HBM holds execution
 and replica regions. Rank-matched Connectors copy KV between hosts;
 each host's DRAM connects to its own Connectors. Blue paths carry
 KV; dashed paths carry control.}
 \label{fig:system-overview}
\end{figure}


\subsection{Tracker}
\label{sec:tracker}

% 中文：一个 Agent 通过多步完成其任务，每步对应一次模型调用。Tracker 按调用粒度
% 记录执行和历史使用关系，按 Agent 剩余工作的粒度估计其对工作流完成时间的影响。
An Agent completes its assigned task through a sequence of steps, each
corresponding to one model invocation. Tracker records execution and
history use at invocation granularity, and estimates remaining Agent
work to determine workflow urgency.

% 中文：本文考虑 Agent 任务 DAG 预先已知且固定的工作流。Agent 集合及依赖关系已知，
% 但每个 Agent 内部的调用次数和完成时间未知。
We target workflows with a fixed Agent-task DAG known before execution.
Agent identities and dependencies are fixed, while each Agent's
invocation count and completion time are unknown.

\subsubsection{Remaining-Time Prediction}
\label{sec:tracker-prediction}

% 中文：以 (w,a,r) 标识工作流 w 中 Agent a 的第 r 次调用，当依赖满足后，workflow runtime 将这次 invocation 提交给 serving 系统调度执行。
% 后续调用可以复用兼容前缀的已选择 KV；调用结束不等于历史失效，直到 Agent 结束。
Tracker identifies invocation $r$ of Agent $a$ in workflow $w$ as
$v=(w,a,r)$. Once its dependencies are satisfied, the workflow runtime
submits the invocation to the serving system for scheduling. Successive
invocations of the same Agent may reuse selected KV for a compatible
history prefix. Completing an invocation does not end the Agent's
history; Tracker tracks that history until the Agent finishes its task.

% 中文：每次调用提交时，Tracker 根据本次调用的信息预测 Agent 距离完成任务还需多久。
% 一次调用只预测一次，下一次调用到来时再更新，直到 Agent 结束。
At each invocation submission, Tracker predicts how long the Agent will
take to finish its task using the invocation's input and the Agent's
observed history. Each invocation triggers one prediction, which is
updated at each subsequent invocation until the Agent finishes its task.

% 中文：将预测建模为基于上下文的回归问题。文本编码器提取任务描述、本次新增输入
% 和近期已接受输出的语义表示，再与提交时已知的模型、角色、输入长度、已完成调用
% 次数、历史调用的长度和耗时，以及 Agent 已运行时间结合。
% 回归器直接预测 Agent 的剩余完成时间，供后续工作流 Slack 计算使用。
We formulate this prediction as context-conditioned regression. A text
encoder represents the Agent's task description, newly appended input,
and recent accepted outputs. A regressor combines this semantic
representation with features available at submission: the model and
Agent role, input length, completed-call count, previous calls' lengths
and timings, and elapsed Agent time. It directly predicts Agent $a$'s
remaining completion time, $\widehat R_{a,r}$, for workflow slack
computation.

% 中文：训练样本将每次调用提交时的特征与该次提交到 Agent 任务结束的实际时间配对。
Training examples pair these submission-time features with the observed
duration from submission of invocation $r$ to completion of Agent $a$'s
task.

% 中文：尚未启动的 Agent 使用同模型、同角色历史任务的时长先验 D。
% 首次调用提交时，以该先验初始化剩余时间估计 R，并记录提交时间 tau。
% 每次调用触发的预测返回后，以新预测替换 R，tau 始终取该预测对应的调用提交时间。
For an unstarted Agent, Tracker uses a duration prior $\bar D_a$ from
historical tasks with the same model and role. Its first invocation
submission initializes $\widehat R_a$ to this prior and $\tau_a$ to the
submission time. Each prediction $\widehat R_{a,r}$ replaces
$\widehat R_a$ when available, with $\tau_a$ set to invocation $r$'s
submission time.

\subsubsection{Slack Computation}
\label{sec:tracker-slack}

% 中文：注册工作流时，runtime 从固定编排图导出所有 Agent 任务的工作流 ID、Agent ID
% 和前驱 Agent ID 列表；边 a→b 表示 b 等待 a 完成任务后才能开始。
% runtime 报告任务完成事件，Tracker 移除已完成节点，得到剩余 DAG。
% 该 DAG 每个节点是一个尚未完成的 Agent 任务，其内部可以有多次 invocation。
At workflow registration, the runtime exports each Agent task's workflow
ID, Agent ID, and predecessor Agent IDs from the fixed orchestration graph.
An edge $a\to b$ means that $b$ waits for Agent $a$ to finish its task.
The runtime reports task completion, allowing Tracker to remove finished
tasks and maintain the residual DAG $G_t=(A_t,E_t)$. Each node in this
graph is one unfinished Agent task, which may contain multiple invocations.

% 中文：算法以剩余 Agent DAG、时长先验、已启动 Agent 的最新估计及提交时间、误差裕量为输入。
% 未启动 Agent 保留完整时长先验；已启动 Agent 的估计换算到当前时刻。
% 正向计算最早完成时间，反向计算最晚完成时间。
% 后向遍历以预测的工作流完成时间为终点；每个 Agent 的剩余工作只计算一次。
% pred 和 succ 分别表示剩余前驱与后继；g 是经校准的 Slack 高估误差裕量。
Algorithm~\ref{alg:slack-estimation} computes guarded Agent slack using
forward and backward critical-path passes~\cite{kelley1959cpm}.
Unstarted Agents retain their full duration priors, while estimates for
started Agents are adjusted to the current time $t$. The algorithm counts
each Agent's remaining work once and uses the predicted workflow finish
$T$ as the backward-pass endpoint. The two passes compute earliest and
latest finishes, $EF_a$ and $LF_a$. Here $\operatorname{pred}(a)$ and
$\operatorname{succ}(a)$ denote residual predecessors and successors,
and $g_a$ is a calibrated upper quantile of slack overestimation.

% 中文：所需时长估计缺失、过期或无效时，整个工作流暂不授予正 Slack。
% 正向遍历：EF = 本节点剩余时间 + 最晚完成的前驱；后向遍历：
% LF = 后继最晚开始时间的最小值，汇点从 T 开始回推。
% 最终 Slack 为 max(0, LF - EF - g)；两次遍历的复杂度均为 O(|A_t|+|E_t|)。
\par\medskip\noindent\begin{minipage}{\columnwidth}
\raggedright
\refstepcounter{designalgorithm}\label{alg:slack-estimation}
\textbf{Algorithm \thedesignalgorithm: Compute Agent slack}
\par\smallskip
\begin{tabularx}{\columnwidth}{@{}r@{\quad}X@{}}
\toprule
 & \textbf{Input:} $G_t=(A_t,E_t)$; Agent states; priors $\bar D_a$;
 estimates $(\widehat R_a,\tau_a)$; guards $g_a$; time $t$ \\
 & \textbf{Output:} guarded slack $\boldsymbol{s}=\{s_a:a\in A_t\}$ \\
\midrule
1 & $s_a\gets0,\quad \forall a\in A_t$ \\
2 & \textbf{for} $a\in A_t$ \textbf{do} \\
3 & \quad \textbf{if} $a$ is unstarted \textbf{then} $d_a\gets\bar D_a$ \\
4 & \quad \textbf{else} $d_a\gets\widehat R_a-(t-\tau_a)$ \\
5 & \textbf{if} any required estimate is missing or invalid, or
 any $d_a\le0$, \textbf{then return} $\boldsymbol{s}$ \\
6 & $\pi\gets\textsc{TopologicalOrder}(G_t)$ \\
7 & \textbf{for} $a$ in $\pi$ \textbf{do} \\
8 & \quad $EF_a\gets d_a+
 \max\bigl(\{0\}\cup\{EF_p:p\in\operatorname{pred}(a)\}\bigr)$ \\
9 & $T\gets\max\bigl(\{0\}\cup\{EF_a:a\in A_t\}\bigr)$ \\
10 & \textbf{for} $a$ in reverse $\pi$ \textbf{do} \\
11 & \quad $LF_a\gets
 \min\bigl(\{T\}\cup\{LF_b-d_b:b\in\operatorname{succ}(a)\}\bigr)$ \\
12 & \quad $s_a\gets\max\{0,LF_a-EF_a-g_a\}$ \\
13 & \textbf{return} $\boldsymbol{s}$ \\
\bottomrule
\end{tabularx}
\end{minipage}\par\medskip

% 中文：调用继承所属 Agent 的 Slack。
Invocation $v=(w,a,r)$ inherits its Agent's slack, $s_v=s_a$.

% 中文：预测和 Slack 更新在 Controller 中异步执行，调用调度和 token 生成无需等待。
Prediction and slack updates run asynchronously in the controller, so
invocation dispatch and token generation do not wait for them.

\subsection{Router and Connector}
\label{sec:router}

% 中文：Router 根据队列负载和前缀 KV 的位置调度调用，与此同时 Router 还会决定如果 Failure 发生时，该 Agent 调用 retry 到哪个DP Replica
% 这一选择也决定了如果要备份该调用对应的 KV，备份在哪里。
% retry 目标位于执行节点的受保护故障域之外，选择时考虑负载、可用内存和复制路径竞争。
Router schedules invocations based on queue load and prefix locality.
It also selects a retry DP replica for each Agent invocation in case of
failure. If the invocation's KV is selected for backup, it is placed at
this retry target. The target lies outside the execution worker's
protected failure domain and is chosen based on load, available memory,
and copy-path contention.

% 中文：已选 KV 存于 retry 目标的 GPU 内存（Peer HBM）或主机内存（Peer DRAM）。
% 配置的前 L 层中的已选区域按 KV Master 的优先级优先放入 HBM；溢出区域和后续层
% 的已选区域放入 DRAM。两层都是部分备份，每个已选区域只有一个稳定备份。
% 容量不足时，Router 将缺口返回 KV Master，重新选择保护区域。
At this target, selected KV resides in GPU memory (\emph{Peer HBM}) or
host memory (\emph{Peer DRAM}). Selected regions in the configured first
$L$ layers prefer Peer HBM according to KV Master's priorities; overflow
and selected regions in later layers use Peer DRAM. Both tiers hold
partial backups, with one stable backup per selected region.
If the selected set does not fit, Router returns the capacity shortfall
to KV Master for reselection.

% 中文：各 TP rank 的 Connector 执行 KV 复制并向 KV Master 报告完成状态。
% 故障后 Router 使旧执行失效，并将调用重定向到 retry 目标；Engine 直接复用目标
% HBM 中的 KV，Connector 将 DRAM 中的 KV 搬入 GPU。目标成为执行节点后，
% Router 为后续备份重新选择 Peer。
Rank-local Connectors execute KV copies and report completion to KV
Master. After a failure, Router fences the old execution and redirects
the invocation to the retry target. The engine reuses resident Peer HBM
KV, while Connectors move DRAM-resident KV into GPU memory.
If that target becomes the execution worker, Router selects a new peer
for subsequent backups.

% 中文：上述路由与传输机制处理已选状态；接下来由 KV Master 决定保护哪些
% Agent 历史，以及保留其中哪些 KV 区域。
These routing and transfer mechanisms operate on selected state.
KV Master determines which Agent histories and which KV regions to
protect, as described next.

\subsection{KV Master}
\label{sec:master}

% 中文：KV Master 根据工作流收益分配有限的备份内存：先选择要保护的 Agent 历史，
% 再选择其中值得保留的 KV 区域，分别对应 O1 和 O2。它使用 Tracker 的剩余时间与
% Slack，以及 Router 的 retry 目标评估方案，并管理副本的更新、保留和回收。
KV Master allocates backup memory according to its predicted benefit to
workflow completion. It first selects Agent histories (O1), then KV
regions within those histories (O2). Tracker supplies remaining-time
estimates and slack, while Router supplies retry destinations. Master
evaluates the resulting protection plans and manages replica lifetimes.

\subsubsection{Agent Selection}
\label{sec:agent-protection}

% 中文：候选来自尚未完成任务的 Agent 所产生的 KV 历史。先比较历史恢复耗时和
% Tracker 给出的 Slack，优先检查可能暴露在工作流完成路径上的恢复延迟。
% 同一故障可能影响汇合点的多个前驱，因此同时考虑单个 Agent 和需要共同保护的小组。
Master starts from materialized histories of unfinished Agents.
It compares the profiled restoration delay with Tracker's guarded slack
$s_a$ to prioritize histories whose recovery may delay the workflow
(Equation~\ref{eq:isolated_penalty}). A failure can interrupt several
predecessors of the same join, so candidates include both individual
histories and small groups of Agents whose recovery delays must be
reduced together.

% 中文：最终选择比较完整方案对工作流完成时间的影响。P 包含已选区域及其放置位置；
% F_w^0 是正常执行时的预测完成时间，F_w^f 是源 worker f 故障后的预测完成时间。
% 目标函数同时计入正常运行和故障恢复，用相同的工作流、Router 目标和权重比较候选方案。
For a protection plan $P$, including selected regions and their
locations, let $F_w^0(P)$ be workflow $w$'s predicted finish during normal
execution and $F_w^f(P)$ its predicted finish after source worker $f$
fails. Master compares plans using
\begin{equation}
 J(P)=(1-\lambda)\sum_w F_w^0(P)
       +\lambda\sum_f\pi_f\sum_w F_w^f(P),
 \label{eq:objective}
\end{equation}
% 中文（续）：lambda 控制恢复目标的权重；pi_f 是预先设定并归一化的故障场景权重，
% 默认等权，不将其解释为实测故障概率。
where $\lambda\in[0,1]$ controls the recovery weight and $\pi_f$ are fixed,
normalized failure-scenario weights, equal by default. Candidate
comparisons use the same workflow snapshot, Router targets, and weights;
the weights express policy preferences rather than measured failure
probabilities.

% 中文：评估以 Tracker 的剩余 Agent 工作为基准，每个 Agent 只计一次，再计入方案
% 改变的 KV 复制、读取和重建开销。故障评估同时考虑该故障影响的调用，以及它们在
% retry 目标实际共享的队列和带宽；就绪调用采用固定顺序，以 Slack 和请求 ID 打破平局。
% 汇合点仍等待所有前驱完成，因此方案收益不能简单相加各 Agent 的独立恢复收益。
The evaluation uses Tracker's remaining Agent work as its baseline,
counting each Agent once and accounting for changes in KV copy,
retrieval, and reconstruction costs. Each failure scenario includes the
affected invocations and the queues and bandwidth they share at their
retry targets. A fixed ready-time order uses slack and request ID to
break ties. Join completion remains the maximum predecessor finish,
so the value of a group is evaluated together.

% 中文：正常复用收益取决于 Router 选定执行位置的实际 KV 获取成本，原执行 KV 已可用时
% 不重复计算备份收益。候选 Agent 进入下一步 KV 区域选择，待放置方案形成后评估
% J(P)-J(P')。若时间预测不可用，则以估计的调用耗时节省/字节排序，并保留容量约束。
Normal reuse contributes value when a replica reduces the KV retrieval
or reconstruction cost of a submitted invocation; state already
available in execution HBM earns no additional backup benefit.
Agent candidates proceed to region selection and placement before
Master evaluates their marginal benefit, $J(P)-J(P')$.
If timing estimates are unavailable, Master orders candidates by
profiled invocation time saved per byte under the same capacity rules.

\subsubsection{KV Region Selection}
\label{sec:region-selection}

% 中文：KV 区域对应某一层的历史 token 块，以及所需的 TP 分片和注意力元数据。
% 对每个候选 Agent，Master 先从已完成物化的区域中形成部分备份集合 S_a，
% 再交由 Router 放置；选择 Agent 不代表为其分配完整检查点。
A KV region identifies one layer and one historical token block,
together with its required TP shards and attention metadata.
For each candidate Agent $a$, Master forms a selected set $S_a$ from
sealed, materialized regions before asking Router to place it.
Selecting an Agent therefore opens its region candidates rather than
allocating a complete checkpoint.

% 中文：先用近期已完成的原生注意力访问记录生成候选。区域的访问因子等于选中该区域
% 的选择单元数，除以该区域在因果关系上可见的选择单元数。
Master uses recent completed native-attention observations to identify
frequently accessed regions. For region $b$ in layer $\ell$, its access
factor is
\begin{equation}
 \widehat p_{a,\ell,b}=
 \frac{\text{eligible selection units that selected }b}
      {\text{selection units for which }b\text{ was visible}}.
 \label{eq:demand}
\end{equation}
% 中文（续）：计数遵循模型原生选择粒度，跨 head 重复访问按单元去重，在剪枝前确定
% 因果可见性。样本不足时使用层级先验；稠密注意力中可见块的访问因子为 1。
% 访问因子用于排列候选，最终保留价值由恢复耗时评估决定。
Counts follow the model's native selection unit, with head duplicates
counted once and causal visibility determined before pruning.
Small denominators use a calibrated layer prior; visible blocks in
dense attention have access factor one. The factor orders candidates,
while their retention value comes from the restoration-time evaluation.

% 中文：候选包括单个区域、逐步扩大的热点集合，以及恢复这些区域所需的依赖组合。
% 连续的历史需求窗口保留访问相关性和需求揭示顺序；统计无效时用完整因果需求范围
% 进行保守评估，但仍由收益和容量决定备份哪些区域。
Master considers individual regions, progressively larger hot sets,
and dependency bundles needed for useful restoration. Recent consecutive
demand windows preserve correlated accesses and their reveal order.
When sparse statistics are invalid, profiling uses the full causal
demand envelope; the selected backup set remains subject to benefit
and capacity constraints.

% 中文：放置后的每个方案区分本地执行 KV、Peer HBM、Peer DRAM 和无存活副本四种情况。
% DRAM 需要传输，未备份状态需要从持久化 token 和可复用状态重建；重建要计入前序层、
% 前序 token、原生选择器及位置元数据等依赖，并对共享重建工作只计一次。
For each proposed placement, the profiler distinguishes an execution-KV
hit, a Peer HBM hit, a Peer DRAM transfer, and missing state.
Missing state is reconstructed from durable tokens and reusable KV,
including required lower-layer activations, preceding KV, and native
selector and positional state. Shared reconstruction work is counted
once in the recovery schedule.

% 中文：恢复评估按照需求揭示、传输和计算的实际先后关系构建时间线，并计入共享资源竞争。
% q_l 表示该层准备好消费历史的时间，a_l 表示所需 KV 全部就绪的时间，c_l 是剩余计算。
% 只有计算不能覆盖的等待 e_l 会推迟该层完成；[x]^+ 表示 max(0,x)。
The profile schedules demand revelation, transfers, reconstruction, and
remaining computation on the retry resources, including their shared
queues and bandwidth. Let $q_\ell$ be the time a layer is ready to consume
history, $a_\ell(P)$ the time its required KV becomes available, and
$c_\ell$ its remaining computation. Exposed waiting and layer completion
are
\begin{equation}
 \begin{aligned}
 e_\ell(P)&=[a_\ell(P)-q_\ell]^+,\\
 f_\ell(P)&=q_\ell+e_\ell(P)+c_\ell,
 \end{aligned}
 \label{eq:overlap}
\end{equation}
% 中文（续）：这里 [x]^+ = max(0,x)。这些等待和完成时间用于计算工作流目标 J；
% HBM 相比 DRAM 的放置优先级取决于每字节能够进一步减少的暴露等待。
where $[x]^+=\max(0,x)$. These timings feed the workflow objective in
Equation~\ref{eq:objective}. For HBM-eligible regions, Master supplies
Router with priorities based on the additional exposed waiting avoided
over DRAM per byte.

% 中文：只有执行引擎能够正确组合部分 KV 与缺失状态重建时，才计入相应的重建节省。
% 若不支持这种恢复方式，则按完整前缀 Prefill 评估，不把不可用的部分副本计为收益。
% 图中展示的是先选择区域、再放入 HBM/DRAM 的构造示例。
The profiler credits partial restoration only when the engine supports
combining retained regions with reconstruction. Otherwise it models
full-prefix Prefill and gives unusable partial state no avoided-compute
credit. Figure~\ref{fig:selective-recovery} illustrates selection followed
by partial placement in the two memory tiers.

% 中文（图示）：12 个逻辑区域中选出 5 个，再按前 L 层限制和容量放入 Peer HBM 或
% Peer DRAM；其余区域没有受管理的备份。这是放置示意，不是性能测量。
\begin{figure}[t]
 \centering
 \includegraphics[width=\columnwidth]{figures/selective-recovery.pdf}
 \caption{Partial selection before placement, on one selected peer, with $L=2$.
 Five of twelve logical regions are selected. Peer HBM holds L1/B0 and
 L1/B2; L2/B1 falls back to Peer DRAM because the illustrative HBM quota is
 full. Selected L3/B0 and L3/B2 also go to Peer DRAM. Unselected cells have
 no managed backup. A demand uses a surviving copy or reconstructs its
 missing dependencies. Each cell denotes all required TP shards;
 this is a constructed state, not a timing result.}
 \label{fig:selective-recovery}
 \Description{A selected-region mask followed by an exclusive Peer HBM/Peer DRAM
 placement map. Two selected first-layer regions reside in Peer HBM, an
 overflowing second-layer region resides in Peer DRAM, and two third-layer
 regions reside in Peer DRAM. Seven regions are unselected. The two tiers
 together store five regions, not a complete twelve-region checkpoint.}
\end{figure}

\subsubsection{Replica Management}
\label{sec:master-online}
\label{sec:execution}

% 中文：新调用、调用完成、Tracker 估计更新及内存预算变化触发重新规划。
% 历史增长按配置的、与块边界对齐的 token 间隔 Z 批量处理，只复制新增或变化的已选区域。
Master revisits selection on invocation arrivals and completions,
Tracker updates, and memory-budget changes. A configured block-aligned
token interval $Z$ batches updates as histories grow.
Only newly selected or changed regions require copying; existing
selected replicas remain available across invocations.

% 中文：候选方案同时包含新增区域、放置位置和为腾出容量需要移除的副本。
% 优先尝试移除单位释放字节造成工作流损失较小的副本，再重新评估完整替换方案。
% 自主替换要求收益超过阈值 epsilon；强制容量收缩可以移除区域，并保留重建路径。
A candidate plan includes its selected regions, destinations, and any
victims needed to fit. Master considers victims in increasing
recomputed workflow loss per released byte, then evaluates the complete
replacement. Optional changes require a benefit greater than a
configured margin $\epsilon$. Mandatory capacity relief may instead
deselect regions, whose missing state remains reconstructible.

% 中文：算法先生成 Agent 候选，再选择区域并请求 Router 放置，重新计算完整方案收益。
% 可行且收益超过阈值的候选按收益/变化区域字节排序；同一物理区域只计一次。
% 所有 TP rank 的 HBM、实际主机的共享 DRAM 以及迁移期间的临时副本都需满足容量限制。
Algorithm~\ref{alg:protection} performs this bounded search.
For candidate $c$, let $\Delta_c=J(P)-J(P_c)$ and let $B_c$ be the physical
bytes of distinct KV regions added, removed, or moved by the candidate.
Feasible candidates with $\Delta_c>\epsilon$ are ordered by
$\Delta_c/B_c$. Feasibility accounts for every TP rank's HBM limit,
shared host DRAM, and transient allocations during migration.

% 中文（算法）：对当前历史、Slack、副本和预算取快照，先处理必要的容量收缩。
% 在规划预算内循环完成 Agent 选择、区域选择、放置及淘汰比较；先过滤低收益候选，
% 再按单位字节收益选择。执行前检查快照和预留，复制完成后发布，使用结束后释放旧副本。
\par\medskip\noindent\begin{minipage}{\columnwidth}
\raggedright
\refstepcounter{designalgorithm}\label{alg:protection}
\textbf{Algorithm \thedesignalgorithm: Update selected replicas}
\par\smallskip
\begin{tabularx}{\columnwidth}{@{}r@{\quad}X@{}}
\toprule
1 & Snapshot histories, slack, replicas, and budgets \\
2 & $P\gets$ current plan with required capacity relief \\
3 & \textbf{while} planning work budget remains \textbf{do} \\
4 & \quad $A\gets\textsc{AgentCandidates}(P)$, including groups \\
5 & \quad $\{S_c\}\gets\textsc{SelectRegions}(A,P)$ \\
6 & \quad For each $c$, form $P_c\gets\textsc{Router.Place}(S_c)$ \\
7 & \quad Fit victims; re-place and evaluate each whole plan \\
8 & \quad Keep feasible candidates with $\Delta_c>\epsilon$ \\
9 & \quad \textbf{if} none remain \textbf{then break} \\
10 & \quad $P\gets P_{c^*}$, where $c^*=\arg\max_c\Delta_c/B_c$ \\
11 & \textbf{if not} \textsc{ReserveIfCurrent}$(P)$ \textbf{then} \\
12 & \quad Requeue update; \textbf{return} \\
13 & Copy newly selected or moved regions to their destinations \\
14 & Publish completed regions; release retired copies after use \\
\bottomrule
\end{tabularx}
\end{minipage}\par\medskip

% 中文：执行前重新检查历史、目录和预算版本，并预留复制所需资源。目录记录层、
% token 区间、TP 分片、兼容性和位置；所有必需分片及元数据就绪后才发布副本。
% 迁移期间新旧副本同时计账，目的副本完成且旧副本的读取或传输结束后才释放旧副本。
Before execution, Master revalidates the history, manifest, and budget
versions and reserves the required resources. The replica manifest
records each region's layer, token range, TP shards, compatibility, and
location. A replica becomes visible only after all required shards and
metadata are ready. During migration, both allocations count until the
destination commits and reads or transfers using the old copy finish.

% 中文：一次调用结束不会自动回收其 Agent 历史中的副本；后续调用可以复用已保留的
% 前缀，并只扩展新增的已选状态。保留集合随使用者和收益变化更新；没有已登记使用者时，
% 可在配置的空闲窗口内保留，内存压力仍可触发替换。工作流结束释放逻辑引用，
% 无剩余使用者的副本在实际读取和传输结束后回收。
Completing an invocation does not retire its Agent's selected replicas.
Later invocations can reuse retained prefixes and extend the selected
state. Master updates retention as consumers and marginal benefits
change; histories without registered consumers may remain for a
configured idle window, subject to memory pressure. Workflow completion
releases logical references, and unreferenced replicas are reclaimed
after active reads and transfers finish.

% 中文：后续调用或 retry 使用目录中已完成的副本，Connector 负责传输，引擎重建缺失
% 状态并按需重放已接受的 token，避免重复输出。目录显式保留未备份的区域，
% 只有经验证完整覆盖所有层及分片的前缀才按完整前缀使用。
For reuse or retry, Master grants access to committed regions,
Connectors transfer them, and the engine reconstructs missing state.
Accepted tokens are replayed where needed without emitting them twice,
following checkpoint-and-replay recovery~\cite{strati2024dejavu,cao2026lumen}.
The manifest preserves gaps in coverage; only a verified all-layer,
all-shard prefix is treated as complete.

\subsection{Adapter}
\label{sec:capacity}

Adapter realizes O3 by determining the HBM replication capacity
foreground execution can support. For takeover rank $k$,
$B^{\rm pool}_k$ excludes weights and fixed allocations,
$B^{\rm serve}_k$ covers the serving high watermark and execution
reservations, and $B^{\rm reserve}_k$ reserves growth and workspace.
The ceiling is
\begin{equation}
 C_k=\max\{0,\min(B^{\rm cap}_k,
 B^{\rm pool}_k-B^{\rm serve}_k-B^{\rm reserve}_k)\}.
 \label{eq:ceiling}
\end{equation}
HBM replica use is not part of $B^{\rm serve}_k$ and is not subtracted twice.
Committed, reserved, copying, and retiring allocations all count toward
physical occupancy. Distinct TP ranks have distinct ceilings.

Adapter uses additive increase and multiplicative
decrease~\cite{chiu1989aimd}. Allocation stalls, reserve violations, or
memory-related admission waits multiply quota $B_k$ by $0<\gamma<1$.
After $m$ stable windows below a lower wait threshold, Adapter adds
$\Delta$ bytes if Master reports beneficial early-layer candidates
currently in Peer DRAM or deferred for HBM capacity. The signal is computed
before quota filtering, allowing growth from zero. Updates are
page-aligned and clipped to $C_k$; missing telemetry disables growth.
Adapter changes total quota, while the configured depth $L$ remains
unchanged.

When quota shrinks, Master stops new HBM admissions and asks Router
to move the least valuable HBM placements to the same peer's DRAM if
their regions remain selected. The plan reserves that host's DRAM
space; Connectors copy and validate the destination. Master then
redirects the manifest and releases the HBM allocation after readers
drain. If that peer's DRAM cannot accommodate the migration, Master
reselects the finite backup set, explicitly removing regions and their
predicted protection benefit as needed.
Their token histories still permit reconstruction. Mandatory pressure
relief need not improve $J$, but cannot bypass active physical leases.
Foreground admission uses actual occupancy; a reduced quota alone
does not free memory. Pending copies may be cancelled, and allocations
wait when immediate physical space is unavailable.

The four components specify a design. Integrated measurements are
still needed to establish recovery gains, cross-invocation reuse, and
normal-service overhead.
